\documentclass[12pt,a4paper]{article} %12 кегль, формат бумаги А4, тип статья

\usepackage[T2A]{fontenc} %поддержка кириллицы
\usepackage[utf8]{inputenc} %кодировка символов юникод
\usepackage[english,russian]{babel} %список рабочих языков
\usepackage{cmap} %для поиска в pdf
\usepackage[left=2cm,right=2cm,top=2cm,bottom=2cm]{geometry} %размеры отступов от краев страницы
\usepackage{graphicx} %возможность работать с графикой
\usepackage{enumitem} %для работы со списками
\usepackage{float} %для точного размещения таблиц и графики
\usepackage{caption} %для создания подписей таблиц и графики
\usepackage{amsmath,amssymb,amsthm,amsfonts} %штуки для математики
\usepackage{extarrows} %чтобы подписи над и под нестандартными стрелками писать короче в коде
\usepackage{makeidx} %собирает ключевые слова для предметного указателя
\usepackage[hidelinks]{hyperref} %убирает выделение ссылок; добавляет гиперссылки в оглавление, перекрестные ссылки и URL
\usepackage{setspace} %устанавливает отступы
\linespread{1.5} %длина межстрочного интервала
\vspace{125mm} %длина абзацного отступа
\usepackage{titlesec}
\titleformat{\section}{\normalfont\Large\bfseries\centering}{}{0pt}{} % эти две строчки убирают циферки перед названиями разделов в основном тексте и выравнивают названия секций по центру
\usepackage{paratype} %чтобы точно был поиск и шрифт хороший

\author{Автор конспекта: Краснов Борис} %автор
\title{Виденский И.В.: введение в теорию пространств Харди} %название
\date{Осень 2026} %дата создания документа

\newtheorem{theorem}{Теорема}
\newtheorem{lemma}{Лемма}
\newtheorem{definition}{Определение}
\newtheorem{proposition}{Предложение}
\newtheorem{consequence}{Следствие}
\newcommand{\R}{\mathbb{R}}
\newcommand{\Co}{\mathbb{C}}
\newcommand{\E}{\mathbb{E}}
\newcommand{\D}{\mathbb{D}}
\newcommand{\T}{\mathbb{T}}
\newcommand{\PR}{\mathbb{P}}
\newcommand{\charf}[2][t]{\varphi_{#2}(#1)}
\newcommand{\F}{\mathcal{F}}
\renewcommand{\H}{\mathcal{H}}
\newcommand{\M}{\mathfrak{M}}

\begin{document}
	\maketitle
	\tableofcontents
	\clearpage
	
	\section{Введение}
	\subsection{Основные обозначения и определения}
	$$\D:=\{z\in\Co\mid|z|<1\}\quad\T:=\{z\in\Co\mid|z|=1\}$$
	За $\H(\Omega)$ будем обозначать аналитичные функции в области $\Omega\in\Co$. Если $f\in\H(\D)$, то положим $f_r(z):=f(rz)$
	\begin{definition}
		Пространство Харди с показателем $p\in(0,+\infty)$ является множество аналитичных функций в круге с нормой:
		$$H^p:=\{f\in\H(\D)\mid||f||_{H^p}:=\sup\limits_{0<r<1}||f||_{L^p(\T,r)}<+\infty\}$$
		$$||f||_{L^p(\T,r)}:=\oint\limits_{|z|=r}|f(w)|^pm(dw),\quad m(dw):=\frac{\lambda}{2\pi}$$
		Также $H^{\infty}:=\{f\in\H(\D)\mid||f||_{L^{+\infty},\D}<+\infty\}$
	\end{definition}
	Известно, что $H^p$ при $1\leq p\leq+\infty$ банахово (в следствии банаховости $L^p$), а также $H^{\infty}$ является алгеброй. Также из вложенности пространств Лебега с конечной мерой имеем: $H^{p_1}\subset H^{p_2}$ при $p_1\geq p_2$.
	\subsection{Классическая теория (1920-1940)}
	Основоположники: Харди, Литтльвуд, Сего, М. и Ф. Рисс, В.И. Смирнов.
	\subsubsection{Граничные значения}
	Пусть $\zeta\in\T$, тогда за $S(\zeta)$ обозначим угол Штольца
	$$S(\zeta):=\{z\in\D\mid|z-\zeta|<k(1-|z|),\quad k>1\}$$
	Будет показано, что для функций из пространств Харди предел значений в угле Штольца, при приблежении к вершины есть значение в вершине п.в. Так называемая теорема о некасательном пределе.
	\subsubsection{Множество нулей}
	Мы хотим по заданной последовательности точек в единичном круге найти аналитичную функцию, которая зануляется в них. Тогда оказывается, что для существования такой функции необходимо и достаточно, чтобы последовательность нулей $Z=\{z_n\}_{n=1}^{+\infty}$ удовлетворяла условию Бляшке:
	$$(B)\quad \sum_{n=1}^{+\infty}(1-|z_n|^2)<+\infty\quad (Z\in(B))$$
	А в качестве подходящей функции можно взять произведение Бляшке:
	$$B^Z(z):=\prod_{n=1}^{+\infty}\Phi_{z_n}(z)\quad\Phi_{a}(z)=\frac{a-z}{1-\overline{a}z}\frac{\overline{a}}{|a|}$$
	Последний множитель в преобразовании Мебиуса нужен для удобства, чтобы в нуле получать значение $|a|$.
	\subsubsection{Факторизация}
	Функцию из пространства Харди можно представить в виде произведения трех компонент: ее нулей - произведение Бляшке, внешней функции, которая по модулю п.в. равна модулю самой функции на $\T$, а также некой сингулярной части.
	\subsection{Эра Карлесона}
	Ну человек придумал какую-то теорию связанную с одноименными мерами:
	$$\sup\limits_{0<r<1}\sup\limits_{\zeta\in\T}\frac{\mu(\{z\mid|z-\zeta|<r\})}{r}<+\infty\Leftrightarrow H^p\subset L^p(\mu),\,p\in(0,+\infty)$$
	$\mu$ - некая борелевская конечная мера в единичном шаре.
	
	Также он решил вопрос интерполяции, а именно если мы возьмем функцию из $H^{\infty}$ и произвольную последовательность точек в круге, тогда ее образ будет лежать в $l_{\infty}$, то есть получаем вложение $R(H^{\infty})\subset l_{\infty}$, где
	$$R(f):=\{f(z_n)\}_{n=1}^{+\infty}$$
	
	Следующий логичный вопрос а есть ли равенство, то есть для любой ли последовательности из $l_{\infty}$ найдется функция из $H^{\infty}$, которая отображает $\{z_n\}_{n=1}^{\infty}$ в данную.
	
	Следующий вопрос это некfя проблема короны и как бы противна мне не была алгебра, но все же слова писать проще чем формулы в техе, так что пусть нам дана конечная последовательность функций из $H^{\infty}$ с нормой ограниченной единицей, тогда возможно ли, что идеал порожденный этими функциями есть все пространство и каковы нижние оценки норм для этого у исходных функций.
	\section{Востановление гармонических функций по граничным значениям}
	Пусть $u(x,y)\in\R$ гармоническая функция в единичном круге. Тогда известно (надо решить простой диффур), что существует гармонически сопряженная к ней в единичном круге $v(x,y)$ и вместе они образуют аналитичную функцию в круге\\ $f(z)=u(x,y)+iv(x,y)$.
	
	Введем обозначение:
	$$h^p:=\{u(x,y)\in C^2(\D)\mid\triangle u(x,y)=0,\,||u||_{h^p}:=\sup\limits_{0<r<1}||u(re^{i\theta})||_{L^p(\T,m)}<+\infty\},\,p\in(0,+\infty]$$
	Тогда оказывается, что принадлежность функции классу Харди равносильна принадлежности ее вещественной и мнимой части соответсвующим $h^p$.
	
	Для начала научимся восстанавливать функцию по граничным значениям и начнем с простого случая:
	\begin{theorem}
		Пусть $u(x,y)\in\R$ гармоническая в некотором круге радиуса $R>1$. Тогда она представляется в виде свертки:
		$$u(re^{i\theta})=(P_r*u)(\theta),\quad r\in(0,1)$$
	\end{theorem}
	\begin{proof}
		Итак, у нас есть гармонически сопряженная функция, то есть исходная является вещественной частью некой аналитичной функции в расширенном круге. Тогда можно представить аналитичную функцию в виде ряда и взять его вещственную часть:
		$$f(z)=u(x,y)+iv(x,y)=\sum_{n=0}^{+\infty}a_nz^n$$
		$$u(re^{i\theta})=Re\,f(re^{i\theta})=\sum_{n=-\infty}^{\infty}A_nr^{|n|}e^{in\theta}$$
		$$
		A_n=
		\begin{cases}
		\frac{a_n}{2},\,n>0 \\
		\frac{a_0}{2},\,n=0 \\
		\frac{\overline{a_n}}{2},\,n<0
		\end{cases}
		$$
		Теперь воспользуемся, что область аналитичности шире единичного круга и тогда имеем представление
		$$A_n=\int_{-\pi}^{\pi}u(e^{it})e^{-int}m(dt)$$
		Тогда пользуясь равномерной и абсолютной сходимостью получаем:
		$$u(re^{i\theta})=\int_{-\pi}^{\pi}u(e^{it})\left(\sum_{n=-\infty}^{+\infty}e^{in(\theta-t)}r^{|n|}\right)m(dt)=(P_r*u)(\theta)$$
		И не сложно посчитав геометрические прогрессии получаем:
		$$P_r(\theta)=\frac{1-r^2}{1+r^2-2r\cos(\theta)}$$
	\end{proof}
	\begin{proposition}
		Система ядер Пуассона $\{P_r\}_{0<r<1}$ образует аппроксимативную единицу, то есть
		\begin{enumerate}
			\item $P_r(\theta)>0$
			\item $\displaystyle\int_{-\pi}^{\pi}P_r(t)m(dt)=1$
			\item При $0<\delta<\pi$ имеем $\lim\limits_{r\rightarrow1-}\sup\limits_{\{\pi\geq|\theta|\geq\delta\}}|P_r(\theta)|=0$
		\end{enumerate}
	\end{proposition}
	\begin{proof}
		Очень очевидно.
	\end{proof}
	Замечание: то есть при $r\rightarrow1-$ ядро Пуассона выражается в дельта-меру.
	\begin{definition}
		Пусть $f\in L^1(\T,m)$, тогда оператор Пуассона:
		$$u(re^{i\theta}):=\int_{-\pi}^{\pi}P_r(\theta-t)f(t)m(dt)$$
		Пусть $\mu\in M_{\Co}(\T)$ - множество борелевских конечных комплексных зарядов, тогда оператор Пуассона:
		$$u(re^{i\theta}):=\int_{-\pi}^{\pi}P_r(\theta-t)\mu(dt)$$
	\end{definition}
\end{document}